%\documentclass{Math_Note}
%\setlength{\parindent}{2em}

%\title{\textbf{Mathmatical Analysis Zorich Comprehensions Chapter 2}}
%\author{Buce-Ithon}
%\newdateformat{mydate}{\twodigit{\THEDAY}{ }\shortmonthname[\THEMONTH], \THEYEAR}
%\date{\today}

%\begin{document}
% Title page
%\maketitle

% Content
%\newpage
%\tableofcontents
%\newpage

% \setcounter{section}{-1}
\newpage
\section{Chapter2. The Real Numbers}
\subsection{Comprehensions}
The main content covers the basic properties of real numbers, the definitions and properties of 
rational and irrational numbers, natural numbers and mathematical induction, and, most importantly, 
several equivalence principles describing the completeness of real numbers (more generally, the density of 
real numbers). (Zorich calls these fundamental lemmas, which serve as the foundation for constructing 
real number theory, but calling them principles is fine.) Finally, some information on Cantor's set theory is added.

The method of describing the completeness of real numbers using Didekind partitioning, which we learned in 
Fehkingoltz's "Calculus Tutorial (I)" can also be incorporated into this.
\newline\newline
\fbox{\begin{minipage}{0.95\textwidth}
\textbf{Basic Lemmas of the Real Number Completeness Theorem}
\begin{enumerate}
    \item[1.] (The least upper (lower) bound principle) Every non-empty set of real numbers that is bounded from above 
    (below) has a unique upper (lower) bound.
    \item[2.] (Cauchy-Cantor Principle) (The nested interval lemma) The set of closed intervals on the real number axis 
    whose length $\to 0$ determines a unique point (a common point that belongs to all these closed intervals).
    \item[3.] (Borel-Lebesgue Principle) (The finite covering lemma) Every system of open intervals covering a closed 
    interval contains a finite subcovering.
    \item[4.] (The Bolzano-Weierstrass Principle) (The limit point lemma) Every bounded infinite set of real numbers has 
    at least one limit point.
\end{enumerate}
\begin{nt*}
1. If the length of the closed intervals is not approaching $0$, there are still existing common points, but they are 
not unique. \\
2. If infinite is limit point, then the Bolzano-Weierstrass Principle is suit for unbounded infinite set of real numbers.
\end{nt*}
\end{minipage}}
\newline\newline
2 important properties of natural numbers $\mathbb{Q}$.
\newline\newline
\fbox{\begin{minipage}{0.95\textwidth}
\textbf{2 Important Properties of Natural Numbers}
\begin{enumerate}
    \item[Thm1.] (The Foundamental Theorem of Arithmetic) $\forall n\in\mathbb{N}$, it could be represented uniquely as: 
    $n=p_1^{\alpha_1}\cdots p_k^{\alpha_k}$, where $p_j$ are different prime numbers, $\alpha_j$ are positive integers.
    \item[Thm2.] (The Principle of Archimedes) Given $h>0$, then $\forall x\in\mathbb{R}$, $\exists$ unique $k\in\mathbb{Z}$, s.t. 
    $(k-1)h\leq x\leq kh$.
\end{enumerate}
\end{minipage}}
\newline\newline
Cantor Set Theory.
\newline\newline
\fbox{\begin{minipage}{0.95\textwidth}
\textbf{Countable and Uncountable Sets}
\begin{enumerate}
    \item[Proposition 1.] card(Countable Set)=card($\mathbb{N}$)=$\aleph_0$, card(Uncountable Set)=card($\mathbb{R}$)=$\aleph_1$ 
    (cardinality of the continuum).
    \item[Proposition 2.] Any subset of a at most countable set (card($X$)$\leq$card($\mathbb{N}$)) is at most countable; the at most 
    countable union of at most countable sets is at most countable.
    \item[Proposition 3.] (Corollary) The direct sum (or direct product) of countable sets is countable. (Corollary) The set of 
    natural numbers or algebraic numbers is countable.
    \item[Thm 1.] (Cantor's Theorem) card($\mathbb{N}$) < card($\mathbb{R}$).
    \item[Proposition 4.] (Corollary) the irrational numbers or transcendental numbers exisst, and the set of all of them is uncountable.
\end{enumerate}
\end{minipage}}

\newpage
\subsection{Problems Sets}
\fbox{\begin{minipage}{0.95\textwidth}
\textbf{An interesting example-the Cantor set}
\begin{ex*}
    Let $x\in[0,1]\subset\mathbb{R}$, its ternary expansion is $x=0.a_1a_2a_3\cdots\, (a_i\in\{0,1,2\})$ satisfying: 
    $\forall i\in\mathbb{N},a\neq 1$, then the set of all such elements is called the Cantor set.
\end{ex*}
\end{minipage}}
\newline\newline
\fbox{\begin{minipage}{0.95\textwidth}
\begin{nt*}
    The Cantor set is uncountable (It's easy: the elements of it can be onto mapped to all binary decimals on $[0,1]$), 
    and we'll then learn: Cantor et is maesure zero in sense of Lebesgue.
\end{nt*}
\end{minipage}}

%\end{document}